\subsection{Identidades representadas, perfiles y declaraciones}\label{sec:diseno-participantes-tipificados}

El directorio separa tres identidades: la cuenta autora que actúa en Qadra, la persona u órgano representado dentro de un expediente y la ficha que le atribuye un rol declarado. Una identidad representada puede tener varias fichas de roles diferentes en ese expediente, pero no concede acceso al sistema. Sus revisiones son independientes de las de cada ficha. La referencia de una ficha conserva UUID, revisión y huella exactos de la identidad empleada; consultar la identidad actual no sustituye retrospectivamente esa referencia.

Una persona física registra nombre conocido o una etiqueta de persona no identificada con motivo, CURP conocida o desconocida con motivo y soporte documental exacto. Un órgano institucional registra denominación, identificador declarado o motivo de desconocimiento y soporte; no se le inventa una CURP. El localizador de página o sección es una declaración legible del operador. Contacto y protección documentados se vinculan a documentos cifrados, sin copiar su contenido al índice ni construir un canal de comunicación. La política documental existente gobierna su consulta; no se promete reserva frente a un rol que ya tiene permiso para leer esos soportes.

El catálogo contempla once perfiles: imputado, víctima u ofendido, defensor, Ministerio Público, asesor jurídico, juez de control, tribunal de enjuiciamiento, perito, policía, supervisión de medidas cautelares y otro. Cada variante define campos propios y estados de desconocimiento admisibles. La validación comprueba forma, límites y compatibilidad de esos valores; no consulta registros civiles o profesionales ni acredita una designación judicial.

Antes de aceptar una identidad nueva o reemplazada, el sistema presenta candidatos por nombre conocido, identificador declarado, huella de certificado y pareja de huella documental y localizador. Los homónimos no se fusionan. Reutilizar una identidad exige consultar su detalle y revisión; declararla distinta exige motivo y soporte para todos los demás candidatos. La revisión abarca el directorio completo mediante una huella de sus cabezas, aunque la pantalla solo muestre una página. Se rechaza una ficha adicional del mismo tipo de rol para la misma identidad, incluidas las archivadas. Estos controles detectan coincidencias registradas, sin demostrar identidad civil ni ausencia de duplicados en la realidad.

La ficha manual conserva su historia original. Completarla agrega una revisión tipificada, sin fabricar revisión cero, reasignar una identidad ya vinculada o retornar a valores manuales. La secuencia revisión, preparación y confirmación mantiene las expectativas y el contenido exacto revisado. Un conflicto conserva el borrador y requiere decisión explícita; una respuesta incierta se concilia con la revisión exacta y su evidencia, sin repetir automáticamente ni deducir éxito por igualdad de nombres.

Defensor y juez de control requieren una firma personal sobre una declaración canónica del prototipo. El titular selecciona un certificado público admitido por la CA interna, exporta los bytes de la declaración y firma fuera de Qadra; el servidor recibe únicamente el certificado y la firma separada. Esta evidencia comprueba control criptográfico de la clave bajo el perfil interno y fija los valores y revisiones declarados. No equivale a FIREL ni acredita por sí misma identidad, profesión, nombramiento o representación. El criterio específico de FIREL de la \cref{tab:rf-06} conserva su alcance pendiente; la credencial interna no se presenta como su sustituto.

{\raggedright Las decisiones se precisan en \rutarepo{docs/adr/0025-case-subjects-and-typed-participants.md} y \rutarepo{docs/adr/0026-internal-participant-declarations.md}. La implementación y sus límites se desarrollan en la \cref{sec:impl-participantes-tipificados}.\par}
